The recent evolution of network technologies has led to the rise of low-latency applications such as cloud gaming, telemedicine, cloud robotics, and cloud virtual reality (VR). Cloud VR eliminates the need for powerful local processing hardware by offloading computation to remote servers, enabling the use of lightweight and cost-effective VR headsets. This architecture opens avenues for diverse innovations, including immersive gaming experiences, enhanced virtual collaboration, and advanced training simulations across various domains. Cloud VR has particularly gained significant importance in recent years \cite{nextmscVirtualReality} due to its ability to deliver immersive, high-resolution VR content to end-users. Cloud VR applications, however, come with stringent technical requirements. They typically require high-resolution content streaming (up to 4K-8K) with high bandwidth demands ($\geq 80$ Mbps) and ultra-low latency ($\leq 20$ ms) to ensure a seamless and interactive user experience. Meeting these requirements is essential to avoid issues like lag, reduced interactivity, and user discomfort, which can culminate in cybersickness \cite{huawei}.  

These demanding requirements pose significant challenges for Cloud VR applications when delivered over wireless networks, particularly Wi-Fi networks. Although recent Wi-Fi standards (e.g., Wi-Fi 6) enable higher bandwidth and lower latency under ideal line-of-sight conditions, real-world performance is hindered by several limitations. These include bandwidth variability, jitter, additional latency, and packet loss, which are exacerbated by environmental factors such as 1) signal attenuation that occurs due to obstacles that absorb the Wi-Fi signal (e.g., walls, furniture); 2) interference that may arise from nearby access points (APs) or coexisting IoT devices using Wi-Fi or Bluetooth protocols, particularly in dense urban environments; or 3) network congestion that can result from high traffic loads on the same AP or simultaneous access by multiple users. These factors degrade the \ac{QoE} for end-users, causing delays, black-edge events, or poor synchronization between movements and visuals, all of which undermine the immersive nature of VR applications \cite{li2020performance, lee2024adaptive, rossi2024subjective}

Given the growing importance of Cloud VR applications, it is imperative to detect and diagnose  performance degradation issues to ensure consistent and reliable sessions. Accurate root-cause diagnosis can help improve network architectures to support emerging low-latency applications or enable \ac{ISP}) to develop intelligent APs capable of diagnosing and mitigating impairments dynamically. Traditionally, root-cause diagnosis has been performed using expert-based rules \cite{hochenbaum2017automatic, chen2019matrix}. In this approach, network experts analyze \ac{KPIs} using predefined heuristics, thresholds, or conditions to identify anomalies and determine their root causes. While effective in simple scenarios, this method is manual, time-consuming, and does not scale well to modern, complex network environments where vast amounts of KPIs are generated continuously.  

In recent years, advancements in machine learning (ML) have revolutionized root-cause diagnosis by automating the analysis of large-scale KPI data. Machine learning models can automatically process time-series KPI data to perform root-cause diagnostic efficiently. \ac{TSC} techniques, in particular, have emerged as effective tools for this task, as they can capture temporal patterns and relationships essential for identifying each cause. Several categories of TSC techniques have been proposed, including distance-based (e.g., DTW), kernel-based (e.g., SVM), shapelet-based approaches, tree-based algorithms or deep learning models leveraging MLPs, CNNs, RNNs, and GNNs \cite{ismail2019deep, faouzi2022time,mohammadi2024deep}. However, supervised TSC methods rely on annotated datasets, which are often expensive and time-consuming to generate, as they require domain expertise to label vast amounts of KPI data. To address the challenges of annotation, researchers are increasingly shifting toward \ac{SSL} for unsupervised time-series classification. SSL has achieved remarkable performance in fields like computer vision and \ac{NLP} by learning meaningful representations of data without relying on labeled samples. Unlike supervised methods, self-supervised approaches define pretext tasks (e.g., contrastive learning or self-prediction) to extract underlying features from raw data. One of the most popular SSL frameworks is \ac{CL}, which uses data augmentation to generate multiple views from a single instance. Positive pairs (augmented views of the same instance) are brought closer together in the feature space, while negative pairs (views of different instances) are pushed apart. Contrastive learning has demonstrated significant success in representation learning for time-series data and forms the foundation of many state-of-the-art unsupervised TSC methods.  

Inspired by these advancements, we propose \ac{RAID}, a two-stage machine learning framework for root-cause diagnosis of performance degradation in Cloud VR applications over Wi-Fi networks: in the first stage, an anomaly detection model uses contrastive learning to differentiate between normal and anomalous time-series KPIs. In the second stage, a lightweight supervised classifier that distinguishes between different types of anomalies detected in Stage 1. We empirically demonstrate that this two-stage framework outperforms existing competing methods, even when trained with limited labeled data. The evaluations are conducted on KPI datasets collected from the real-world Cloud VR testbed, presented on Chapter \ref{chap:measurements}, featuring a Cloud VR application built on NVIDIA CloudXR architecture. The system streams VR content over a Wi-Fi connection using a commercial access point and an Oculus Quest 2 headset. Our proposed solution offers a scalable and efficient method for root-cause diagnosis, paving the way for improved network support for low-latency applications like Cloud VR.  

\contributions{
Specifically, the key contributions of this chapter are as follows:
\begin{itemize}
    % \item We design a controlled Wi-Fi testbed that replicates real-world network impairments, enabling the study of their impact on Cloud VR performance.
    \item We introduce a novel two-stage framework that combines contrastive learning-based anomaly detection with supervised classification to effectively detect and diagnose Wi-Fi impairments in Cloud VR environments.
    \item We perform extensive empirical evaluations using time series KPI datasets collected from our testbed that replicates real-world network impairments, comparing our proposed solution with state-of-the-art time series classification models.
    \item Our solution demonstrates strong performance even in low-label scenarios, highlighting its ability to generalize with minimal supervision. Additionally, it offers a balanced trade-off between moderate training time and low inference latency, making it well-suited for real-time deployment in practical applications.
\end{itemize}
}
